
\section[Teoremas de comparación cohomológica]{Teoremas de comparación cohomológica y teoremas de existencia}
\label{XII.4}
\marginpar{326}

El objeto de este número es volver a demostrar los resultados de
\cite[\no 2 th\ptbl 5 y th\ptbl 6]{XII.3}; estos generalizan al caso de un
esquema propio los teoremas establecidos en \cite[\no 12]{XII.10} cuando
$X$ es proyectivo, y los extienden al caso relativo. En
\cite[ch\ptbl VIII \no 3]{XII.7} se demuestran resultados más generales
relativos a esquemas relativos propios sobre un espacio analítico.

Recordemos que la cohomología de {\v{C}}ech utilizada en
\cite[\no 12]{XII.10} coincide con la cohomología usual tanto en el caso
algebraico como en el analítico (EGA~III 1.4.1 y
\cite[II 5.10]{XII.5}).
\subsection{}
\label{XII.4.1}
\setcounter{proposition}{1}% The first theorem in this source section is XII.4.2.
Sea $f\colon X\to Y$ un morfismo de~$\CC$-esquemas localmente de tipo
finito y consideremos el diagrama conmutativo
$$
\xymatrix{ X^\an \ar[r]^\varphi \ar[d]_{f^\an} & X \ar[d]^f
\\ Y^\an \ar[r]^\psi & Y \quoi.}
$$
Si $F$ es un $\cal{O}_X$-módulo, para todo entero $p\geq0$ se tienen
morfismos
$$
\R^p f_* F \lto{i} \R^p f_* (\varphi_* F^\an )
\lto{j} \R^p (f\cdot\varphi)_* F^\an \lto{k} \psi_*
(\R^p f_*^\an F^\an ) \quoi,
$$
donde $i$ se deduce del morfismo canónico $F\to\varphi_*F^\an$, y $j,k$
son ``homomorfismos de borde'' de sucesiones espectrales de Leray. Al
compuesto $k.j.i$ se le asocia un morfismo canónico
\begin{equation*}
\label{eq:XII.4.1.1}
\tag{4.1.1} {\theta_p \colon (\R^p f_* F)^\an \to \R^p f_*^\an
(F^\an )}
\end{equation*}

\begin{theoreme}
\label{XII.4.2}
Sean
\marginpar{327}
$f\colon X\to Y$ un morfismo propio de~$\CC$-esquemas localmente de tipo
finito y $F$ un $\cal{O}_X$-módulo coherente. Entonces, para todo entero
$p\geq0$, el morfismo~\eqref{eq:XII.4.1.1}
$$
\theta_p \colon (\R^p f_* F)^\an \to \R^p f_*^\an (F^\an)
$$
es un isomorfismo.
\end{theoreme}

1) \emph{Caso en que $f$ es proyectivo}. La demostración es análoga a la de
\cite[\no 13]{XII.10}. Recordémosla brevemente. Nos reducimos al caso en
que $X$ es un espacio proyectivo
\ifthenelse{\boolean{orig}}
{tipo $\PP^r_Y$}
{de tipo $\PP^r_Y$}%
~sobre~$Y$. Sean $\cal{Y}=Y^\an$ y $\cal{P}=\PP^r_\cal{Y}$; se demuestra
primero que
$$
f_*^\an \cal{O}_\cal{P} = \cal{O}_\cal{Y} \quoi, \qquad \R^p f_*^\an
(\cal{O}_\cal{P}) = 0 \qquad \text{para } p>0
$$
En efecto, para verificar las relaciones anteriores podemos reducirnos al
caso en que $\cal{Y}$ es una bola $\cal{B}$ de un espacio afín
$\cal{E}^n$. Consideremos el ``recubrimiento estándar'' $\{\cal{U}_i\}$
de~$\cal{P}$ por $r+1$ abiertos isomorfos a
$\cal{B}\times\cal{E}^r$. Como estos abiertos son de Stein, para todo
entero $p\geq0$ se tienen isomorfismos
$$
\H^p (\{ \cal{U}_i \}, \cal{O}_\cal{P}) \isomto \H^p (\cal{P},
\cal{O}_\cal{P})
$$
Podemos expresar entonces las secciones del haz estructural
$\cal{O}_\cal{P}$ sobre los abiertos $\cal{U}_i$ y sobre sus intersecciones
en términos de series de Laurent; un cálculo sencillo demuestra que
$$
\H^0 (\cal{P}, \cal{O}_\cal{P} ) \isomto \H^0 (\cal{Y},
\cal{O}_\cal{Y} ) \quoi, \qquad \H^p (\cal{P}, \cal{O}_\cal{P} ) = 0
\qquad \text{para } p>0
$$
La demostración concluye repitiendo \cite[\no 12 lema~5]{XII.10}, con los
grupos de cohomología sustituidos por los haces de cohomología.

2)
\marginpar{328}
\emph{Caso en que $f$ es propio}. Utilizamos EGA~III 3.1.2 para reducirnos
al caso proyectivo. Sea $\cal{K}$ la categoría de los
$\cal{O}_X$-módulos coherentes para los que $\theta_p$ es un isomorfismo
para todo $p\geq0$. Basta demostrar que, para toda sucesión exacta
$0\to F'\to F\to F''\to0$ dos de cuyos términos pertenecen a~$\cal{K}$,
lo mismo ocurre con el tercero; que todo sumando directo de un objeto
de~$\cal{K}$ pertenece a~$\cal{K}$; y que, para todo punto $x$ de~$X$,
puede encontrarse un objeto $F$ de~$\cal{K}$ tal que $F_x\neq0$.

\ifthenelse{\boolean{orig}}
{El primera condición}
{La primera condición}%
~resulta de aplicar el lema de los cinco al siguiente diagrama conmutativo,
cuyas filas son exactas:
$$
\xymatrix@=.7cm{ \ar[r] & (\R^p f_* F')^\an \ar[d] \ar[r] &(\R^p f_*
F)^\an \ar[d] \ar[r] & (\R^p f_* F'')^\an \ar[d] \ar[r] &
(\R^{p+1} f_* F' )^\an \ar[r] \ar[d] &\, \\ \ar[r] & \R^p
f_*^\an F'^\an \ar[r] &\R^p f_*^\an F^\an \ar[r] & \R^p
f_*^\an F''^\an \ar[r] & \R^{p+1} f_*^\an F'^\an \ar[r] &
\quoi,}
$$
y la segunda condición se verifica de manera análoga.

Para verificar la tercera condición, podemos limitarnos al caso en que $X$
es un esquema irreducible con punto genérico~$x$. Podíamos haber supuesto
desde el principio que~$Y$
era noetheriano. Por el lema de Chow (EGA~II 5.6.1), existe un
$Y$-esquema proyectivo~$X'$ y un morfismo proyectivo sobreyectivo
$g\colon X'\to X$. Por otra parte, existe un entero~$n$ tal que
$\R^p g_* (\cal{O}_{X'}(n))=0$ para todo $p>0$ y que el morfismo canónico
$g^*g_*(\cal{O}_{X'}(n))\to\cal{O}_{X'}(n)$ es sobreyectivo
(EGA~III 2.2.1). Si ponemos $F=g_*(\cal{O}_{X'}(n))$, el haz~$F$ posee la
propiedad requerida. En efecto, $F_x\neq0$; además, como la sucesión
espectral de Leray
$$
\R^p f_* (\R^q g_* (\cal{O}_{X'} (n))) \To \R^{p+q}
(\ifthenelse{\boolean{orig}}{f.g}{f\cdot g})_* (\cal{O}_{X'} (n))
$$
es degenerada, se tiene un isomorfismo
$$
\R^p f_* F \isomto \R^p (\ifthenelse{\boolean{orig}}{f.g}{f\cdot g})_* (\cal{O}_{X'} (n))
$$
\marginpar{329}%
Como en el caso algebraico, se tiene un isomorfismo canónico
$$\R^p f_*^\an F^\an \isomto
\R^p (\ifthenelse{\boolean{orig}}{f.g}{f\cdot g})_*^\an (\cal{O}_{X'}(n)^\an ),$$
y el diagrama
$$
\xymatrix@C=1cm{
(\R^p f_* F)^\an \ar[d]_{\theta_p} \ar[r]^-{\sim} & (\R^p (\ifthenelse{\boolean{orig}}{f.g}{f\cdot g})_* (\cal{O}_{X'} (n)))^\an \ar[d]_{\psi_p} \\
\R^p f_*^\an F^\an \ar[r]^-{\sim} & \R^p (\ifthenelse{\boolean{orig}}{f.g}{f\cdot g})_*^\an
(\cal{O}_{X'}(n)^\an)}
$$
es conmutativo. Por~1), $\psi_p$ es un isomorfismo; por tanto, también lo
es~$\theta_p$, lo cual concluye la demostración.

\begin{corollaire}
\label{XII.4.3}
Sean $X$ un $\CC$-esquema propio y $F$ un $\cal{O}_X$-módulo coherente.
Entonces, para todo entero $p\geq0$, el morfismo canónico
$$
\H^p(X,F)\to \H^p(X^{\an},F^{\an})
$$
es un isomorfismo.
\end{corollaire}

\begin{theoreme}
\label{XII.4.4}
Sea $X$ un $\CC$-esquema propio. El funtor que asocia a cada
$\cal{O}_X$-módulo coherente~$F$ su imagen inversa $F^\an$ sobre~$X^\an$
es una equivalencia de categorías.
\end{theoreme}

1) \emph{El funtor es plenamente fiel}. Sean, en efecto, $F$ y $G$ dos
$\cal{O}_X$-módulos coherentes. El morfismo canónico
$$
\Hom_{\cal{O}_X}(F,G)\to \Hom_{\cal{O}_{X^{\an}}}(F^{\an},G^{\an})
$$
se identifica con el morfismo canónico
$$
\H^0(X,\SheafHom_{\cal{O}_X}(F,G))\to
\H^0(X^{\an},
\ifthenelse{\boolean{orig}}
{\SheafHom_{\cal{O}_X}(F,G)}
{\SheafHom_{\cal{O}_{X^\an}}(F^\an,G^\an)})
$$
\marginpar{330}%
(EGA~$0_{\textup{I}}$~6.7.6). Como $\SheafHom_{\cal{O}_X}(F,G)$ es
coherente, de~\Ref{XII.4.3} resulta que este morfismo es biyectivo.

2) \emph{El funtor es esencialmente sobreyectivo}. Cuando $X$ es
proyectivo, la afirmación resulta de \cite[\no 12~th\ptbl 3]{XII.10}. El
caso general se reduce al anterior mediante el lema de Chow
(EGA~II~5.6.1). Sean, en efecto, $X'$ un $\CC$-esquema proyectivo,
$f\colon X'\to X$ un morfismo proyectivo sobreyectivo y $U$ un abierto
denso de~$X$ tal que $f$ induce un isomorfismo $f^{-1}(U)\simeq U$.
Razonamos por inducción noetheriana sobre~$X$; podemos, pues, suponer que,
para todo haz coherente $\cal{G}$ sobre~$X^\an$ para el cual existe un
subconjunto cerrado $Y$ de~$X$, distinto de~$X$, que satisface
$Y^\an\supset\Supp\cal{G}$, existe un haz coherente $G$ sobre~$X$ tal que
$G^\an\simeq\cal{G}$.

Sea $\cal{F}$ un haz coherente de $\cal{O}_{X^\an}$-módulos, y sean
$\cal{K}$ y $\cal{L}$ los haces coherentes definidos por la exactitud de la
sucesión
$$
0\to\cal{K}\to\cal{F}\to f_*^\an f^{\an *}\cal{F}\to\cal{L}\to 0
$$
Como $X'$ es proyectivo, existe un $\cal{O}_{X'}$-módulo coherente $F'$ tal
que $F'^\an\simeq{f^\an}^*\cal{F}$; de~\Ref{XII.4.2} se deduce entonces
un isomorfismo $(f_*F')^\an\simeq f_*^\an{f^\an}^*\cal{F}$. Como
$\cal{K}|U^\an$ y $\cal{L}|U^\an$ son nulos, existen
$\cal{O}_X$-módulos coherentes $K$ y $L$
\ifthenelse{\boolean{orig}}
{tal que}
{tales que}%
~$K^\an\simeq\cal{K}$ y $L^\an\simeq\cal{L}$. Por~1), el morfismo
$f_*^\an{f^\an}^*\cal{F}\to\cal{L}$ procede de un único morfismo
$f_*F'\to L$; sea $I=\Ker(f_*F'\to L)$. El haz $\cal{F}$ es entonces una
extensión de~$I^\an$ por~$K^\an$, y basta ver que esta extensión procede
\marginpar{331}
por imagen inversa de una extensión de~$I$ por~$K$. Basta, por tanto,
demostrar que el morfismo canónico
\begin{equation*}
\label{eq:XII.4.*}
\tag{$*$}
\ifthenelse{\boolean{orig}}
{\Ext_{\cal{O}_X}^q(I,K)^\an}
{\Ext_{\cal{O}_X}^q(I,K)}
\isomto
\Ext_{\cal{O}_{X^{\an}}}^q(I^{\an},K^{\an})\quad
\ifthenelse{\boolean{orig}}{q\neq 1}{q=1}
\end{equation*}
es biyectivo. Ahora bien, se tienen isomorfismos
$\mathbf{Ext}_{\cal{O}_X}^q(I,K)^\an\isomto
\mathbf{Ext}_{\cal{O}_{X^{\an}}}^q(I^{\an},K^{\an})$
para todo entero $q\geq0$ (EGA~$0_{\textup{III}}$~12.3.5), y un morfismo
de sucesiones espectrales
\begingroup
\shorthandoff{>}% Spanish babel otherwise treats the Xy-pic arrow tip as a shorthand.
$$
\xymatrix@R=.7cm@C=.5cm{
\H^p(X,\mathbf{Ext}_{\cal{O}_X}^q(I,K))\ar@{=>}[r]\ar[d] &
\Ext_{\cal{O}_X}^{p+q}(I,K)\ar[d]\\
\H^p(X^{\an},\mathbf{Ext}_{\cal{O}_{X^{\an}}}^q(I^{\an},K^{\an}))\ar@{=>}[r] & \Ext_{\cal{O}_{X^{\an}}}^{p+q}(I^{\an},K^{\an}).
}
$$
\endgroup
Este morfismo es un isomorfismo porque, por~\Ref{XII.4.3}, lo es en los
términos $E^{pq}_2$, y esto demuestra la biyectividad de~$(*)$.

\begin{corollaire}
\label{XII.4.5}
El funtor que asocia a cada $\CC$-esquema propio $X$ el espacio analítico
$X^\an$ es plenamente fiel.
\end{corollaire}

Debemos demostrar que, si $X$ e $Y$ son dos $\CC$-esquemas propios, la
aplicación canónica
$$
\Hom_{\CC}(X,Y)\to\Hom(X^{\an},Y^{\an})
$$
es biyectiva. Ahora bien, dar un morfismo de~$X$ en~$Y$ (\resp de
$X^\an$ en $Y^\an$) equivale a dar su grafo, \ie un subesquema cerrado
$Z$ de~$X\times Y$ (\resp un subespacio analítico cerrado $\othercal{Z}$
de~$X^\an\times Y^\an$), tal que la restricción de la primera proyección
$X\times Y\to X$ a~$Z$ (\resp la de
$X^\an\times Y^\an\to X^\an$ a~$\othercal{Z}$) sea
\marginpar{332}
un isomorfismo. Como dar un subesquema cerrado de~$X\times Y$ (\resp un
subespacio analítico cerrado de~$X^\an\times Y^\an$) equivale a dar un haz
coherente de ideales de~$\cal{O}_{X\times Y}$ (\resp de
$\cal{O}_{X^\an\times Y^\an}$), el corolario resulta de~\Ref{XII.4.4}.

\begin{corollaire}
\label{XII.4.6}
Sea $X$ un $\CC$-esquema propio. El funtor que asocia a cada esquema
finito (\resp finito étale) $X'$ sobre~$X$ el espacio analítico $X'^\an$
es una equivalencia de la categoría de los esquemas finitos (\resp finitos
étale) sobre~$X$ con la categoría de los espacios analíticos finitos
(\resp finitos étale) sobre~$X^\an$.
\end{corollaire}

En efecto, dar un morfismo finito $X'\to X$ (\resp
$X'^\an\to X^\an$) equivale a dar un haz coherente de álgebras sobre
$\cal{O}_X$ (\resp sobre $\cal{O}_{X^\an}$)
\cite[\no 19~\S 5~th\ptbl 2]{XII.4}. El corolario resulta, pues,
de~\Ref{XII.4.4} en el caso sin paréntesis, y el caso entre paréntesis se
deduce de él teniendo en cuenta~\Ref{XII.3.1}~(iii).
